\section*{Chapter \ref{ch:g10:lewis-and-shape} --- Lewis Structures and the Shape of Molecules}

\begin{solution}{exo:g10:lewis-and-shape:1}
A pair of \omterm{def:g9:inside-the-atom:nucleus}{electrons} shared by two \omterm{def:g7:atoms-and-molecules:atom}{atoms}. A pair of \omterm{def:g10:electron-shells:valence}{valence electrons}
belonging to one \omterm{def:g7:atoms-and-molecules:atom}{atom} only, not shared.
\end{solution}

\begin{solution}{exo:g10:lewis-and-shape:2}
H: 1 (one \omterm{def:g9:inside-the-atom:nucleus}{electron} short of its duet). C: 4, N: 3, O: 2 (4, 3 and 2
\omterm{def:g9:inside-the-atom:nucleus}{electrons} short of an octet).
\end{solution}

\begin{solution}{exo:g10:lewis-and-shape:3}
\chemfig{H-\charge{0=\:,90=\:,270=\:}{Cl}}: chlorine has the shared pair and three \omterm{def:g10:lewis-and-shape:lone-pair}{lone
pairs}, 8 \omterm{def:g9:inside-the-atom:nucleus}{electrons}.
\end{solution}

\begin{solution}{exo:g10:lewis-and-shape:4}
$5 + 3 \times 1 = 8$ \omterm{def:g9:inside-the-atom:nucleus}{electrons}, 4 pairs (three bonds, one \omterm{def:g10:lewis-and-shape:lone-pair}{lone pair}).
\end{solution}

\begin{solution}{exo:g10:lewis-and-shape:5}
\ce{CH4}: tetrahedron. \ce{NH3}: pyramid. \ce{H2O}: bent. \ce{CO2}:
linear.
\end{solution}

\begin{solution}{exo:g10:lewis-and-shape:6}
\chemfig{H-\charge{90=\:,270=\:}{S}-H}: like water, four groups around sulfur of
which two \omterm{def:g10:lewis-and-shape:lone-pair}{lone pairs}: bent.
\end{solution}

\begin{solution}{exo:g10:lewis-and-shape:7}
$2 \times 5 = 10$ \omterm{def:g9:inside-the-atom:nucleus}{electrons}, 5 pairs. A single bond and 4 \omterm{def:g10:lewis-and-shape:lone-pair}{lone pairs} leave
each nitrogen with 6 \omterm{def:g9:inside-the-atom:nucleus}{electrons}; turning two \omterm{def:g10:lewis-and-shape:lone-pair}{lone pairs} into shared pairs
gives a \omterm{def:g10:lewis-and-shape:multiple-bond}{triple bond} and one \omterm{def:g10:lewis-and-shape:lone-pair}{lone pair} on each \omterm{def:g7:atoms-and-molecules:atom}{atom}:
\chemfig{\charge{180=\:}{N}~\charge{0=\:}{N}}.
\end{solution}

\begin{solution}{exo:g10:lewis-and-shape:8}
\chemfig{H-C(-[2]H)(-[6]H)-\charge{90=\:,270=\:}{O}-H}: four bonds around carbon,
tetrahedral; around oxygen two bonds and two \omterm{def:g10:lewis-and-shape:lone-pair}{lone pairs}, bent.
\end{solution}

\begin{solution}{exo:g10:lewis-and-shape:9}
In \ce{CO2} carbon has two groups (two \omterm{def:g10:lewis-and-shape:multiple-bond}{double bonds}) and no \omterm{def:g10:lewis-and-shape:lone-pair}{lone pair}: they
point in opposite directions. In \ce{H2O} oxygen has four groups (two
bonds, two \omterm{def:g10:lewis-and-shape:lone-pair}{lone pairs}) in a tetrahedral arrangement: the two bonds make an
angle.
\end{solution}

\begin{solution}{exo:g10:lewis-and-shape:10}
\chemfig{C(-[3]H)(-[5]H)=C(-[1]H)-[7]H}: around each carbon, three
groups (one \omterm{def:g10:lewis-and-shape:multiple-bond}{double bond}, two single bonds): a flat triangle, angles near
\ang{120}.
\end{solution}

\begin{solution}{exo:g10:lewis-and-shape:11}
Carbon in the centre, four single bonds to chlorine \omterm{def:g7:atoms-and-molecules:atom}{atoms} each with three
\omterm{def:g10:lewis-and-shape:lone-pair}{lone pairs}: four groups around carbon, a tetrahedron.
\end{solution}

\begin{solution}{exo:g10:lewis-and-shape:12}
Ethanol, \chemfig{H-C(-[2]H)(-[6]H)-C(-[2]H)(-[6]H)-\charge{90=\:,270=\:}{O}-H}, and
methoxymethane,
\chemfig{H-C(-[2]H)(-[6]H)-\charge{90=\:,270=\:}{O}-C(-[2]H)(-[6]H)-H}.
\end{solution}

\begin{solution}{exo:g10:lewis-and-shape:13}
A \omterm{def:g10:lewis-and-shape:lone-pair}{lone pair} takes up a little more room than a \omterm{def:g10:lewis-and-shape:lone-pair}{bonding pair} and pushes the
bonds closer together. Ammonia has one \omterm{def:g10:lewis-and-shape:lone-pair}{lone pair}, water two: the angle
shrinks from methane to ammonia to water.
\end{solution}

\begin{solution}{exo:g10:lewis-and-shape:14}
\chemfig{H-C~\charge{0=\:}{N}}: two groups around carbon, linear.
\chemfig{H-C(-[6]H)=\charge{45=\:,315=\:}{O}}: three groups around carbon, a flat
triangle.
\end{solution}

\begin{solution}{exo:g10:lewis-and-shape:15}
$3 \times 6 = 18$ \omterm{def:g9:inside-the-atom:nucleus}{electrons}, 9 pairs:
\chemfig{\charge{180=\:,270=\:}{O}=\charge{90=\:}{O}-\charge{0=\:,180=\:,270=\:}{O}}. The central oxygen
has three groups (a \omterm{def:g10:lewis-and-shape:multiple-bond}{double bond}, a single bond, a \omterm{def:g10:lewis-and-shape:lone-pair}{lone pair}): the
\omterm{def:g7:atoms-and-molecules:molecule}{molecule} is bent.
\end{solution}

\begin{solution}{pb:g10:lewis-and-shape:1}
\textbf{1.} \chemfig{H-C(-[2]H)(-[6]H)-H}, \chemfig{H-\charge{90=\:}{N}(-[6]H)-H},
\chemfig{H-\charge{90=\:,270=\:}{O}-H}.

\textbf{2.} Four groups in each.

\textbf{3.} None for methane, one for ammonia, two for water.

\textbf{4.} Carbon forms four bonds; oxygen forms two (its two other
groups are \omterm{def:g10:lewis-and-shape:lone-pair}{lone pairs}, which a kit does not show with sticks).

\textbf{5.} \chemfig{H-C(-[2]H)(-[6]H)-C(-[2]H)(-[6]H)-H}: one stick joins
the two carbon balls, and six hydrogen balls are needed.

\textbf{6.} Tetrahedron, pyramid, bent.

\textbf{7.} \omterm{def:g10:lewis-and-shape:lone-pair}{Lone pairs} take more room than \omterm{def:g10:lewis-and-shape:lone-pair}{bonding pairs} and squeeze the
bonds together; water, with two \omterm{def:g10:lewis-and-shape:lone-pair}{lone pairs}, has the smallest angle.

\textbf{8.} A carbon ball's holes make \ang{109.5}; water's true angle is
\ang{104.5}: an oxygen ball needs its own two holes at \ang{104.5}.

\textbf{9.} \ang{180}: two groups around carbon point in opposite
directions, and the \omterm{def:g7:atoms-and-molecules:molecule}{molecule} is linear.

\textbf{10.} Two corners not on the same edge are opposite corners of one
face: their distance is the face diagonal, $\sqrt{a^2 + a^2} = a\sqrt{2}$.

\textbf{11.} The centre is halfway along a body diagonal:
$a\sqrt{3}/2$.

\textbf{12.} H--H: $\sqrt{2} \approx \qty{1.41}{cm}$; C--H:
$\sqrt{3}/2 \approx \qty{0.87}{cm}$.

\textbf{13.} Each hydrogen is at a corner of the cube and carbon at its
centre: every corner is at the same distance from the centre.

\textbf{14.} The triangle C\,H$_1$\,H$_2$ is isosceles (C\,H$_1$ = C\,H$_2$);
the line from its apex to the midpoint of the base is perpendicular to the
base.

\textbf{15.} $\sin \widehat{\text{M\,C\,H}_1} = \dfrac{\text{M\,H}_1}{\text{C\,H}_1}
= \dfrac{a\sqrt{2}/2}{a\sqrt{3}/2} = \sqrt{2/3} \approx 0.816$.

\textbf{16.} Half-angle $\approx \ang{54.7}$; angle H--C--H $\approx
2 \times 54.7 = \ang{109.5}$.

\textbf{17.} It is the \ang{109.5} of the tetrahedron in the table.

\textbf{18.} The holes must be drilled at \ang{109.5} from one another.
\end{solution}
